\section{Conclusions}

This paper has established a comprehensive modeling framework for the design and optimization of PDMS-based radiative cooling technology, addressing four progressively complex sub-problems.

\textbf{Problem 1} developed a single-layer PDMS emissivity model using the Transfer Matrix Method with Airy summation. The model reveals that PDMS emissivity in the 8--13 $\mu$m atmospheric window increases from approximately 0.05 at sub-micron thicknesses to 0.89 at 50 $\mu$m, approaching saturation ($\bar{\varepsilon}_{8-13} \approx 0.97$) beyond 200 $\mu$m. The spectral features are dominated by Si--O--Si and Si--C vibrational absorption bands at 7.9, 9.3, and 12.5 $\mu$m.

\textbf{Problem 2} evaluated the radiative cooling performance across thicknesses. The net cooling power peaks at 76.5 W/m$^2$ for a 40 $\mu$m PDMS film, representing the optimal balance between high IR emissivity ($\bar{\varepsilon}_{8-13}=0.938$) and low solar absorptance ($\bar{\alpha}_{\text{solar}}=0.023$). Beyond 50 $\mu$m, diminishing returns set in as increased solar absorption outweighs marginal emissivity gains. Our analysis identifies 40--50 $\mu$m as the recommended thickness range for single-layer PDMS radiative coolers.

\textbf{Problem 3} investigated multilayer structures incorporating reflective and dielectric layers. The Ag (100 nm) / SiO$_2$ (0.5 $\mu$m) / PDMS (50 $\mu$m) configuration provides effective solar reflection while maintaining strong IR emission. The SiO$_2$ spacer layer offers modest but measurable performance enhancement through optical impedance matching. Systematic parameter scanning confirmed the robustness of the optimal thickness parameters.

\textbf{Problem 4} conducted a comprehensive evaluation encompassing cost analysis, economic feasibility, and technology readiness. The optimal design achieves a manufacturing cost of approximately \$40.5/m$^2$ with an estimated payback period of 4--6 years. All required fabrication processes are industrially mature (TRL 8--9). The technology is most suitable for arid and semi-arid regions where the atmospheric window transmittance is highest.

The modeling framework developed herein provides a transferable methodology applicable to other radiative cooling material systems. Future work should focus on experimental validation, nanostructured surface engineering for enhanced spectral selectivity, and integration with building-scale energy simulation for deployment optimization.
